\section{Importance-Weighted Skill-Set Evaluation Under Policy Drift}
\label{app:iw}

In the co-evolution framework, rollouts collected at earlier GRPO steps were generated by a behavior policy $\pi_{\theta_t}$ that differs from the current policy $\pi_{\theta_T}$ at decision time $T$.  To evaluate each candidate skill set under the current policy using all historical data, we apply importance weighting---the standard off-policy correction.  This appendix details the token-level formulation for our GRPO setting.

%-----------------------------------------------------------------------------
\subsection{Data-generating process}
\label{app:iw:data}
%-----------------------------------------------------------------------------

At each GRPO step $t = 1, \ldots, T$:

\begin{enumerate}[leftmargin=*]
    \item A batch of prompts is sampled:
    \begin{equation}
        \{x_{t,1}, \ldots, x_{t,B}\}, \qquad x_{t,i} \sim \mathcal{D},
    \end{equation}
    where $B$ is the batch size (e.g.\ $B=32$) and $\mathcal{D}$ is the task distribution.

    \item For each prompt $x_{t,i}$, we generate $G$ rollouts (e.g.\ $G=8$):
    \begin{equation}
        r_{t,i,1},\; r_{t,i,2},\; \ldots,\; r_{t,i,G}.
    \end{equation}

    \item Each rollout is assigned to one of the two competing skill sets:
    \begin{equation}
        S_{t,i,k} \in \{S_1, S_2\}.
    \end{equation}

    \item Each rollout is a token sequence of length $L_{t,i,k}$, sampled autoregressively from the behavior policy:
    \begin{align}
        r_{t,i,k} &= (a_{t,i,k,1},\; a_{t,i,k,2},\; \ldots,\; a_{t,i,k,L_{t,i,k}}), \\
        a_{t,i,k,\ell} &\sim \pi_{\theta_t}(\cdot \mid x_{t,i},\, S_{t,i,k},\, a_{t,i,k,<\ell}), \qquad \ell = 1, \ldots, L_{t,i,k}.
    \end{align}

    \item A sequence-level binary reward is observed:
    \begin{equation}
        R_{t,i,k} = R(x_{t,i},\, r_{t,i,k}) \in \{0, 1\}.
    \end{equation}
\end{enumerate}

So one rollout sample is a tuple $(x_{t,i},\, S_{t,i,k},\, r_{t,i,k},\, R_{t,i,k})$ with indices $i = 1, \ldots, B$ and $k = 1, \ldots, G$, giving $B \times G$ trajectory samples per GRPO step.

%-----------------------------------------------------------------------------
\subsection{Target quantity}
\label{app:iw:target}
%-----------------------------------------------------------------------------

\begin{definition}[Current-policy value of a skill set]
\label{def:value}
At decision time $T$, the value of skill set $S$ under the current policy is
\begin{equation}
\label{eq:iw_value}
    V_T(S) := \mathbb{E}_{x \sim \mathcal{D}}\, \mathbb{E}_{r \sim \pi_{\theta_T}(\cdot \mid x, S)}\bigl[R(x, r)\bigr].
\end{equation}
\end{definition}

The goal is to estimate $V_T(S_1)$ and $V_T(S_2)$ using all historical rollouts within this cycle $\{1,\cdots,T\}$, then select the skill set with higher estimated value.

%-----------------------------------------------------------------------------
\subsection{Notation summary}
\label{app:iw:notation}
%-----------------------------------------------------------------------------

\Cref{tab:iw_notation} collects the key symbols used in this appendix before the derivation.

\begin{table}[ht]
\centering
\caption{Notation for importance-weighted skill-set evaluation.}
\label{tab:iw_notation}
\begin{tabular}{@{}ll@{}}
\toprule
\textbf{Symbol} & \textbf{Meaning} \\
\midrule
$t = 1, \ldots, T$ & GRPO training step \\
$i = 1, \ldots, B$ & Prompt index within a batch (e.g.\ $B=32$) \\
$k = 1, \ldots, G$ & Rollout index per prompt (e.g.\ $G=8$) \\
$\ell = 1, \ldots, L_{t,i,k}$ & Token position within rollout \\
$x_{t,i}$ & Prompt \\
$S_{t,i,k} \in \{S_1, S_2\}$ & Skill set assigned to rollout \\
$a_{t,i,k,\ell}$ & Token $\ell$ of rollout $(t,i,k)$ \\
$r_{t,i,k} = (a_{t,i,k,1}, \ldots, a_{t,i,k,L})$ & Full rollout token sequence \\
$R_{t,i,k} \in \{0,1\}$ & Sequence-level reward \\
$\pi_{\theta_t}$ & Behavior policy at step $t$ \\
$\pi_{\theta_T}$ & Current (target) policy at decision time \\
$\rho_{t,T,\ell}^{(i,k)}$ & Per-token importance ratio (\cref{eq:token_ratio}) \\
$w_{t,T}^{(i,k)}$ & Sequence-level importance weight (\cref{eq:seq_weight}) \\
$V_T(S)$ & Current-policy value of skill set $S$ (\cref{eq:iw_value}) \\
\bottomrule
\end{tabular}
\end{table}

%-----------------------------------------------------------------------------
\subsection{Importance-weight identity}
\label{app:iw:identity}
%-----------------------------------------------------------------------------

For a rollout $r$ collected at step $t$ under skill set $S_{t,i,k} = S$ and behavior policy $\pi_{\theta_t}$, the standard importance-sampling identity gives
\begin{equation}
    \mathbb{E}_{r \sim \pi_{\theta_T}(\cdot \mid x, S)}\bigl[R(x, r)\bigr]
    = \mathbb{E}_{r \sim \pi_{\theta_t}(\cdot \mid x, S)}\!\left[\frac{\pi_{\theta_T}(r \mid x, S)}{\pi_{\theta_t}(r \mid x, S)}\, R(x, r)\right].
\end{equation}

The sequence-level importance ratio is
\begin{equation}
\label{eq:iw_seq_ratio}
    w_{t,T}^{(i,k)} := \frac{\pi_{\theta_T}(r_{t,i,k} \mid x_{t,i},\, S_{t,i,k})}{\pi_{\theta_t}(r_{t,i,k} \mid x_{t,i},\, S_{t,i,k})}.
\end{equation}

%-----------------------------------------------------------------------------
\paragraph{Token-level decomposition.}
\label{app:iw:token}
%-----------------------------------------------------------------------------

Since language model generation factorizes autoregressively,
\begin{equation}
    \pi_\theta(r \mid x, S) = \prod_{\ell=1}^{L} \pi_\theta(a_\ell \mid x, S, a_{<\ell}),
\end{equation}
the importance ratio decomposes into a product of per-token ratios.

\begin{definition}[Per-token importance ratio]
\label{def:token_ratio}
For rollout $(t,i,k)$ and token position $\ell$, define
\begin{equation}
\label{eq:token_ratio}
    \rho_{t,T,\ell}^{(i,k)} := \frac{\pi_{\theta_T}(a_{t,i,k,\ell} \mid x_{t,i},\, S_{t,i,k},\, a_{t,i,k,<\ell})}{\pi_{\theta_t}(a_{t,i,k,\ell} \mid x_{t,i},\, S_{t,i,k},\, a_{t,i,k,<\ell})}.
\end{equation}
\end{definition}

Then the sequence-level importance weight is
\begin{equation}
\label{eq:seq_weight}
    w_{t,T}^{(i,k)} = \prod_{\ell=1}^{L_{t,i,k}} \rho_{t,T,\ell}^{(i,k)}.
\end{equation}

\paragraph{Log-ratio form.}
In practice, one works in log-space to avoid numerical overflow:
\begin{equation}
\label{eq:log_ratio}
    \log w_{t,T}^{(i,k)} = \sum_{\ell=1}^{L_{t,i,k}} \left[\log \pi_{\theta_T}(a_{t,i,k,\ell} \mid x_{t,i},\, S_{t,i,k},\, a_{t,i,k,<\ell}) - \log \pi_{\theta_t}(a_{t,i,k,\ell} \mid x_{t,i},\, S_{t,i,k},\, a_{t,i,k,<\ell})\right].
\end{equation}

Note that using the exact weights $w_{t,T}^{(i,k)}$ to estimate $V_T(S)$ is unbiased in principle: for each rollout with $S_{t,i,k} = S$, the identity $\mathbb{E}_{\pi_{\theta_t}}\!\left[w_{t,T}^{(i,k)}\, R_{t,i,k}\right] = \mathbb{E}_{\pi_{\theta_T}}\!\left[R_{t,i,k}\right]$ holds.  However, the exact product $\prod_\ell \rho_{t,T,\ell}^{(i,k)}$ can become extremely large or small for long sequences, especially when step $t$ is far from $T$, making the exact IS estimator impractical.  This motivates the surrogate objective in the next subsection.

%-----------------------------------------------------------------------------
\subsection{Surrogate objective and estimator}
\label{app:iw:surrogate}
%-----------------------------------------------------------------------------

Following the same philosophy as PPO~\citep{schulman2017proximal}---which replaces the exact policy gradient with a clipped surrogate objective---we replace the exact importance weights with clipped versions and define the surrogate estimator directly.

\paragraph{Clipped importance weights.}
We consider two clipping strategies.

\textit{Per-token clipping.}  Clip each token-level ratio before taking the product:
\begin{equation}
\label{eq:clip_token}
    \tilde{\rho}_{t,T,\ell}^{(i,k)} = \mathrm{clip}\!\left(\rho_{t,T,\ell}^{(i,k)},\; 1 - \epsilon,\; 1 + \epsilon\right),
\end{equation}
then form the clipped sequence-level weight
\begin{equation}
\label{eq:clip_seq}
    \tilde{w}_{t,T}^{(i,k)} = \prod_{\ell=1}^{L_{t,i,k}} \tilde{\rho}_{t,T,\ell}^{(i,k)}.
\end{equation}
This is the direct analogue of the PPO clipped surrogate, applied per-token.


\paragraph{Surrogate value estimator.}
Using the clipped weights $\tilde{w}_{t,T}^{(i,k)}$ from either~\eqref{eq:clip_seq}, we define the surrogate value estimator in the plain IS form, consistent with PPO/GRPO:
\begin{equation}
\label{eq:iw_surrogate}
    \tilde{V}_T(S) := \frac{1}{N_S}\sum_{t=1}^{T} \sum_{i=1}^{B} \sum_{k=1}^{G} \mathbf{1}\{S_{t,i,k} = S\}\; \tilde{w}_{t,T}^{(i,k)}\; R_{t,i,k},
\end{equation}
where $N_S = \sum_{t=1}^{T} \sum_{i=1}^{B} \sum_{k=1}^{G} \mathbf{1}\{S_{t,i,k} = S\}$ is the total number of rollouts assigned to $S$.

This is the quantity we actually compute.  It is a biased surrogate of the exact target $V_T(S)$ (\cref{def:value}), but the bias vanishes when the policy drift within the evaluation window is small: for recent steps where $\theta_t \approx \theta_T$, the per-token ratios satisfy $\rho_{t,T,\ell}^{(i,k)} \approx 1$ and clipping does not activate, so $\tilde{w} \approx w$ and the surrogate coincides with the exact IS estimator.  This parallels how PPO's clipped surrogate matches the true policy gradient near the old policy.

\paragraph{Skill-set selection rule.}
Given the surrogate estimates for both skill sets, define the estimated gap
\begin{equation}
\label{eq:iw_gap}
    \hat{\Delta}_T := \tilde{V}_T(S_1) - \tilde{V}_T(S_2).
\end{equation}
The decision rule is:
\begin{itemize}
    \item choose $S_1$ if $\hat{\Delta}_T > 0$,
    \item choose $S_2$ if $\hat{\Delta}_T < 0$.
\end{itemize}

%-----------------------------------------------------------------------------
\subsection{Efficient computation}
\label{app:iw:efficient}
%-----------------------------------------------------------------------------

Computing $\tilde{w}_{t,T}^{(i,k)}$ requires access to both $\pi_{\theta_t}$ and $\pi_{\theta_T}$.  The key insight is that old model checkpoints need not be stored if sufficient information is saved at rollout time.

\paragraph{At rollout step $t$.} For each rollout $r_{t,i,k}$, store:
\begin{itemize}
    \item prompt $x_{t,i}$, skill assignment $S_{t,i,k}$, generated tokens $a_{t,i,k,1:L_{t,i,k}}$,
    \item per-token log-probabilities under the behavior policy:
    \begin{equation}
        \log \pi_{\theta_t}(a_{t,i,k,\ell} \mid x_{t,i},\, S_{t,i,k},\, a_{t,i,k,<\ell}), \qquad \ell = 1, \ldots, L_{t,i,k},
    \end{equation}
    \item reward $R_{t,i,k}$.
\end{itemize}

\paragraph{At decision step $T$.} Run one forward pass of the current model $\pi_{\theta_T}$ on all stored trajectories to obtain $\log \pi_{\theta_T}(a_{t,i,k,\ell} \mid x_{t,i},\, S_{t,i,k},\, a_{t,i,k,<\ell})$.  Compute each $\log \rho_{t,T,\ell}^{(i,k)}$ via~\eqref{eq:token_ratio}, apply clipping via~\eqref{eq:clip_seq}, and evaluate the surrogate estimator~\eqref{eq:iw_surrogate}.  This avoids reloading any old checkpoint.
